\chapter{The Coefficient-Majorant Carath\'eodory Disk}
\label{chap:coefficient-majorant-caratheodory-disk}

\status{SOH-G019 proved: a canonical positive-real-part disk and normalized coefficient law.}

\section{Purpose}

Chapter~\ref{chap:quarter-disk-zero-exclusion} established the zero-free disk
\[
 |w|\le \frac14
\]
for the even quotient
\[
 \xi\!\left(\frac12+z\right)=F(z^2),
 \qquad
 F(w)=\sum_{n\ge0}a_nw^n,
\]
and used it to close the negative-inversion paired-spectrum frontier.  The present chapter returns to the open real-rootedness problem and extracts a stronger consequence of the already proved strict coefficient positivity
\[
 a_n>0\qquad(n\ge0).
\]
No table of zeta zeros enters the argument.

\section{The radial coefficient-majorant threshold}

Put
\[
 a_0=F(0)=\xi\!\left(\frac12\right)>0.
\]
For real \(r\ge0\),
\[
 F'(r)=\sum_{n\ge1}na_nr^{n-1}>0,
\]
so \(F\) is strictly increasing on the nonnegative axis.  Moreover, since \(a_1>0\),
\[
 F(r)\ge a_0+a_1r\longrightarrow\infty.
\]
Therefore the equation
\[
 F(R_\star)=2F(0)
\]
has exactly one positive solution.

\begin{definition}[Coefficient-majorant threshold]
The unique number \(R_\star>0\) satisfying
\[
 \boxed{F(R_\star)=2F(0)}
\]
is called the SOH-G019 coefficient-majorant threshold.
\end{definition}

The definition is exact.  The numerical value of \(R_\star\) is not an input to the proof.

\section{Positive real part on the whole disk}

For any \(w\in\C\), coefficient positivity gives
\[
 \left|F(w)-a_0\right|
 \le \sum_{n\ge1}a_n|w|^n
 =F(|w|)-a_0.
\]
Consequently, whenever \(|w|<R_\star\),
\begin{align*}
 \Re F(w)
 &\ge a_0-\left|F(w)-a_0\right|\\
 &\ge 2a_0-F(|w|)\\
 &>0.
\end{align*}

The boundary requires one additional equality check.  If \(|w|=R_\star\), then
\[
 \sum_{n\ge1}a_n|w|^n=a_0.
\]
Equality in the triangle inequality for the nonconstant series would require all terms \(a_nw^n\) to have one common argument.  Since \(a_1>0\) and \(a_2>0\), the first two terms can have the same argument only when \(w=R_\star>0\).  At that point
\[
 F(R_\star)=2a_0>0.
\]
At every other boundary point the triangle inequality is strict.

\begin{theorem}[SOH-G019 positive-real-part disk]
For the even xi quotient \(F\),
\[
 \boxed{\Re F(w)>0\qquad(|w|\le R_\star).}
\]
Hence
\[
 \boxed{F(w)\ne0\qquad(|w|\le R_\star).}
\]
\end{theorem}

\begin{corollary}[Centered xi disk]
Since \(w=(s-\tfrac12)^2\),
\[
 \boxed{
 \left|s-\frac12\right|\le\sqrt{R_\star}
 \quad\Longrightarrow\quad
 \xi(s)\ne0.
 }
\]
\end{corollary}

This is a zero-free neighborhood derived from the coefficient theorem itself rather than from a table of zero ordinates.

\section{Canonical coefficient normalization}

The threshold identity is equivalent to
\[
 \sum_{n\ge1}a_nR_\star^n=a_0.
\]
Define
\[
 p_n:=\frac{a_nR_\star^n}{a_0},
 \qquad n\ge1.
\]
Then every \(p_n\) is strictly positive and
\[
 \boxed{\sum_{n\ge1}p_n=1.}
\]
Thus
\[
 \boxed{
 \frac{F(R_\star\zeta)}{F(0)}
 =1+\sum_{n\ge1}p_n\zeta^n.
 }
\]
The scaled coefficient sequence is therefore exactly normalized on the positive integers.

Since all terms are positive and the sum is one,
\[
 0<p_n<1,
\]
which gives the global envelope
\[
 \boxed{
 a_n<F(0)R_\star^{-n},
 \qquad n\ge1.
 }
\]
This bound is uniform in \(n\) and requires no asymptotic estimate.

\section{Numerical regression}

High-precision bisection of the defining equation gives the diagnostic values
\[
 R_\star\approx30.7037329843450643,
\]
\[
 \sqrt{R_\star}\approx5.54109492648746.
\]
Therefore the G018 radius \(1/4\) is strictly inside the G019 coefficient-majorant disk.  The approximate numbers are regression outputs only; all theorem statements above are analytic consequences of coefficient positivity and the definition of \(R_\star\).

\section{Relation to the remaining frontier}

G019 does not prove that \(R_\star\) is the largest actual zero-free radius of \(F\).  It is the canonical threshold at which the simple radial coefficient budget
\[
 F(r)-F(0)\le F(0)
\]
becomes saturated.  The resulting normalized sequence \((p_n)_{n\ge1}\) provides a new exact coefficient-level object for subsequent PF and Jensen analysis.

\gap{Real-rootedness of \(F\), PF$_3$, PF$_\infty$, and the general SOH-G003 real-rootedness criterion remain open.  No Riemann-Hypothesis conclusion is made in this chapter.}
